\section{アドミッタンス制御}
\label{sec:admittance_control}

本研究で使用する制御系には，提案手法を導入する以前から，外力に対する柔軟な応答を実現するための
\textbf{アドミッタンス制御}が組み込まれている．
本節では，このアドミッタンス制御の基本的な考え方と，本研究における位置づけについて説明する．

アドミッタンス制御は，ロボットが環境から受ける力やトルクを入力とし，それに応じた位置あるいは速度の応答を生成する制御手法である．
一般に，外力 $\mathbf{F}_{\mathrm{ext}}$ に対して，仮想的な質量，粘性，剛性を持つ動特性を設定し，次式のような関係で運動を決定する：

\begin{equation}
M \ddot{\mathbf{x}} + D \dot{\mathbf{x}} + K (\mathbf{x} - \mathbf{x}_{\mathrm{ref}})
= \mathbf{F}_{\mathrm{ext}}
\end{equation}

ここで，$\mathbf{x}$ は制御対象の位置，$\mathbf{x}_{\mathrm{ref}}$ は基準位置，
$M, D, K$ はそれぞれ仮想質量，仮想粘性，仮想剛性を表す．
このように，\textbf{力入力に対して位置応答を生成する}点が，位置入力に対して力応答を生成するインピーダンス制御との違いである．

二足歩行ロボットにおいては，足裏接触や外乱によって環境から複雑な力が作用する．
これらに対して位置を厳密に拘束すると，関節トルクの急激な増加や不安定な挙動を引き起こす可能性がある．
そのため，本制御系ではアドミッタンス制御を用いることで，外力に応じて目標位置を微小に調整し，
ロボットの運動に柔軟性を持たせている．

本研究における制御系では，歩行計画器や安定化制御器によって生成された
重心位置や足先位置の目標値が，全身逆運動学へ入力される前段階で，
アドミッタンス制御による補正を受ける構成となっている．
これにより，接触力や外乱が存在する状況においても，
高位制御で生成された目標を破綻させることなく実行可能となっている．

なお，本研究ではアドミッタンス制御自体の設計やパラメータ調整を主題とはせず，
既存の制御要素としてこれを利用する．
以降で述べる安定化制御および重心制御は，
このアドミッタンス制御が動作している状態を前提として構成されている．